\section{Conclusions}
\vspace{-1.5mm}
 This work uncovers the critical vulnerability of generative reward models for complex reasoning with reference answers, where superficial patterns trigger false positive rewards. We propose a data augmentation strategy that substantially mitigates this issue. We conduct comprehensive experiments revealing that mid-sized judges may offer a better accuracy–robustness trade-off than larger ones, chain-of-thought prompting does not reliably improve robustness, and removing the questions from prompts significantly reduces false positives. These findings provide concrete guidelines for deploying robust generative reward models in RLVR.
%This work uncovers a critical vulnerability in generative reward models where superficial patterns trigger false positive rewards. To mitigate this, we propose a data augmentation strategy and demonstrate that mid-sized judges offer a superior accuracy–robustness trade-off. Our extended analysis further reveals that while inference-time techniques like chain-of-thought and majority voting fail to reduce false positives, removing original questions from prompts significantly improves robustness, particularly for larger models. These findings establish practical guidelines for deploying robust generative reward models in RLVR.